\documentclass[11pt]{report}
\input{preambolo-verifica.tex}

\begin{document}
\chapter*{Verifica sulle equazioni di II grado incomplete}
\begin{flushright}
    \textbf{Fila A}\\
    Tempo stimato: 60 minuti.
\end{flushright}

\begin{flushright}
    Classe: 3B \quad Data: Lunedì, 28 Settembre 2026
\end{flushright} 

\vspace{0.8em}
\regoladoro
\begin{paracol}{2}
    
\section*{Risolvi questi esercizi per arrivare fino a 9}
Trova l'incognita.
\begin{enumerate}
    \item $\dfrac{1-\sqrt2}{\sqrt[4]2}\,n^2=0$
    \item $-\sqrt\pi\,x^2-\dfrac{15}{4}=0$
    \item $9-2z^2=-9$
    \item $(2x+3)^2=5x+9$
    \item $\left(\dfrac x2-2\right)\left(\dfrac x2+2\right)=8$
    \item $-\dfrac{y+3}{4} - (3y-2)^2 = 2y - \dfrac{19}{4}$
\end{enumerate}
\switchcolumn

\section*{Prova a risolvere questi esercizi per aggiungere punti}
Questi esercizi sono leggermente diversi da quelli standard visti in classe. Prova a ragionarci (\ul{senza utilizzare la formula di risoluzione delle equazioni complete}).

Trova l'incognita.

\begin{enumerate}
    \item[7.] $-3(x-1)^5=(x-1)^3$
    \item[8.] $x^2+6x+9=7$
    \item[9.] $x\left(\dfrac {x^2}2 +1\right)(x^3-2)^4\left(2x+\frac14\right)=0$
\end{enumerate}
\end{paracol}
\end{document}
